\section{Deployment-Gap and Cross-Engine Results}
\label{sec:deployment_gap_results}

The frozen M4 controller was evaluated in 375 new one-factor stress
evaluations.  Together with 45 reused nominal evaluations, the panel contains
420 observations over the same three blocks, five checkpoints, and three
force targets.  Every stress trace retained direct TD-MPC2 action authority;
the native simulator force remained the outcome signal while only the
learner-visible observation or registered task frame was perturbed.

\begin{table*}[t]
\centering
\caption{Severe deployment-gap endpoints. Compound counts are per 15 evaluations in each fixed block; task, tracking, and sampled-limit safety counts are over 45 evaluations. Peak force uses the native simulator signal.}
\label{tab:deployment_gap_endpoints}
\scriptsize
\begin{tabular}{lrrrrrrr}
\toprule
Condition & B0 comp. & S1 comp. & S2 comp. & Task & Tracking & Safety & Max peak (N) \\
\midrule
Nominal & 11/15 & 7/15 & 6/15 & 42/45 & 24/45 & 43/45 & 17.905 \\
Force noise 0.035~N RMS & 12/15 & 6/15 & 7/15 & 41/45 & 25/45 & 41/45 & 18.679 \\
Force bias $-0.8$~N & 14/15 & 13/15 & 4/15 & 40/45 & 31/45 & 40/45 & 22.337 \\
Force bias $+0.8$~N & 3/15 & 0/15 & 9/15 & 42/45 & 12/45 & 43/45 & 18.031 \\
Delay 50~ms & 11/15 & 12/15 & 5/15 & 34/45 & 28/45 & 35/45 & 26.084 \\
TCP error 0.10~mm & 11/15 & 7/15 & 7/15 & 44/45 & 25/45 & 44/45 & 18.491 \\
Normal error 5$^{\circ}$ & 9/15 & 8/15 & 5/15 & 30/45 & 24/45 & 41/45 & 21.316 \\
\bottomrule
\end{tabular}
\end{table*}


The signed force bias produced the clearest regulation response.  A
$+0.8$~N bias reduced compound pass from 24/45 under nominal observation to
12/45, including 11/15 to 3/15 in B0 and 7/15 to 0/15 in S1.  The corresponding
paired changes were $-53.33$ percentage points in B0 (95\% pointwise interval
$[-66.67,-40.00]$) and $-46.67$ points in S1 ($[-60.00,-33.33]$).  A
$-0.8$~N bias instead yielded 31/45 compound passes, but it also reduced task
and sampled-limit safety to 40/45 and raised the maximum native-force peak to
22.337~N.  The signed intervention therefore exposes observation bias rather
than a uniformly beneficial operating offset.

Delay and frame error affected different parts of the task.  At 50~ms,
task and sampled-limit safety fell from 42/45 and 43/45 to 34/45 and 35/45,
and the maximum peak reached 26.084~N.  A 5-degree registered-normal error
reduced task completion to 30/45 and compound pass to 22/45.  By contrast,
0.10-mm TCP-position error produced 25/45 compound passes versus 24/45
nominal, and 0.035-N RMS force noise produced 25/45.  These one-factor results
identify signed force calibration, latency, and frame alignment as the main
deployment sensitivities in the evaluated simulation blocks; their mixed
effects across B0, S1, and S2 also rule out a single monotone robustness score.

\subsection{Zero-Shot MuJoCo Stress Test}

The independently implemented MuJoCo port was evaluated only after its
open-loop mismatch had been recorded. Table~\ref{tab:cross_engine_transfer}
reports the 45 zero-shot evaluations and their matched PhysX cells. The five
frozen checkpoints achieve task completion, tracking, and sampled-limit safety
in all 30 flat and inclined evaluations. In the cylindrical block, however,
all 15 evaluations fail each component: no episode completes, every trace
exceeds 15~N, and the maximum peak reaches 20.24~N. The matched PhysX block
is also the weakest of its three surfaces, but it retains 12/15 task
successes, 6/15 tracking and compound passes, and 13/15 sampled-limit safety
passes.

\begin{table*}[t]
\centering
\caption{Zero-shot cross-engine stress results on matched checkpoint--block--target keys. Counts are out of 15 evaluations per block.}
\label{tab:cross_engine_transfer}
\scriptsize
\setlength{\tabcolsep}{5pt}
\renewcommand{\arraystretch}{0.95}
\begin{tabular}{@{}llrrrrr@{}}
\toprule
Engine & Block & Task & Tracking & Sampled safety & Compound & Max. peak (N) \\
\midrule
PhysX  & Flat reference       & 15 & 11 & 15 & 11 & 13.09 \\
MuJoCo & Flat reference       & 15 & 15 & 15 & 15 & 13.32 \\
PhysX  & $10^\circ$ incline  & 15 &  7 & 15 &  7 & 12.30 \\
MuJoCo & $10^\circ$ incline  & 15 & 15 & 15 & 15 & 13.20 \\
PhysX  & Cylinder, $r=1.2$~m & 12 &  6 & 13 &  6 & 17.90 \\
MuJoCo & Cylinder, $r=1.2$~m &  0 &  0 &  0 &  0 & 20.24 \\
\bottomrule
\end{tabular}
\end{table*}


The engine change therefore preserves the complete measured outcome on the
two planar geometries and breaks it consistently on weak curvature. Because
the calibration gate retains a bidirectional tangential-position mismatch,
this contrast is a portability stress result rather than a solver-equivalence
claim. It shows that the flat and inclined capability is not confined to one
engine implementation, while curved contact remains dependent on the policy,
geometry mapping, and contact solver as a coupled system.
